\chapter{Pangkat}\label{ch:g8:powers}

Lipatlah selembar kertas menjadi dua sebanyak $10$ kali (kalau kamu
bisa!): tebalnya berlipat dua setiap kali, dan $10$ kali pelipatan
mengalikannya dengan $2^{10} = 1024$. Pangkat adalah singkatan
perkalian berulang; bab ini menyiapkan notasinya dan aturannya, dengan
peran khusus untuk pangkat $10$.

\section{Definisi}

\begin{definition}[Pangkat]\label{def:g8:powers:def}
Untuk sebuah bilangan $a$ dan bilangan bulat $n \geq 1$:
\[
a^n = \underbrace{a \times a \times \dots \times a}_{n
\text{ faktor}},
\]
dibaca ``$a$ pangkat $n$''; $a$ adalah \emph{bilangan pokoknya}, $n$
adalah \emph{eksponennya}\index{eksponen}. Nama khusus: $a^2$ adalah
``$a$ kuadrat'', $a^3$ adalah ``$a$ pangkat tiga''. Menurut kesepakatan:
\[
a^1 = a,
\qquad
a^0 = 1 \ (a \neq 0),
\qquad
a^{-n} = \frac{1}{a^n}.
\]
\end{definition}

\begin{example}\label{ex:g8:powers:def}
$3^4 = 3 \times 3 \times 3 \times 3 = 81$; \quad
$10^3 = 1000$; \quad $5^{-2} = \frac{1}{25} = 0.04$; \quad
$(-2)^3 = -8$ dan $(-2)^4 = +16$
(\cref{ex:g8:negprod:powers}). Hati-hati: $a^n$ \emph{bukan}
$a \times n$: \ $2^5 = 32$, bukan $10$.
\end{example}

\begin{omfigure}
\begin{tikzpicture}[scale=0.9]
  \foreach \n/\x in {0/0, 1/1.4, 2/2.8, 3/4.2, 4/5.6} {
    \pgfmathsetmacro{\h}{0.28*pow(2,\n)}
    \draw[omDef, fill=omDef!20] (\x,0) rectangle (\x+0.9,\h);
    \node[below, font=\small] at (\x+0.45,0) {$2^\n$};
    \pgfmathtruncatemacro{\v}{pow(2,\n)}
    \node[above, font=\small] at (\x+0.45,\h) {$\v$};
  }
\end{tikzpicture}

{\small Pangkat $2$: tiap batangnya dua kali batang sebelumnya.
Pertumbuhan lewat perkalian berulang melesat jauh lebih cepat daripada
pertumbuhan lewat penjumlahan berulang.}
\end{omfigure}

\section{Aturan eksponen}

\begin{theorem}[Aturan eksponen]\label{thm:g8:powers:rules}
Untuk bilangan pokok $a$ yang bukan \omterm{def:g1:counting:zero}{nol} dan \omterm{def:g8:powers:def}{eksponen} bilangan bulat $m$, $n$:
\[
a^m \times a^n = a^{m+n},
\qquad
\frac{a^m}{a^n} = a^{m-n},
\qquad
\left(a^m\right)^n = a^{m \times n},
\qquad
(ab)^n = a^n b^n .
\]
\end{theorem}

\begin{proof}[Bukti dengan membilang faktornya]
$a^m \times a^n$ menjajarkan $m$ faktor $a$ yang disusul $n$ faktor
lagi: seluruhnya $m + n$. $\left(a^m\right)^n$ mengulang satu blok
berisi $m$ faktor sebanyak $n$ kali: $mn$ faktor. $(ab)^n$ memuat $n$
huruf $a$ dan $n$ huruf $b$, yang dapat dikelompokkan ulang. Aturan
\omterm{def:g3:division:remainder}{hasil baginya} datang dari mencoret $n$ faktor dari $m$ faktornya;
dengan kesepakatan $a^0 = 1$ dan $a^{-n} = \frac{1}{a^n}$, aturan itu
tetap benar bahkan ketika $n \geq m$.
\end{proof}

\begin{example}\label{ex:g8:powers:rules}
\[
2^3 \times 2^5 = 2^8 = 256,
\qquad
\frac{7^6}{7^4} = 7^2 = 49,
\qquad
\left(5^2\right)^3 = 5^6,
\qquad
2^4 \times 5^4 = 10^4 .
\]
Sebuah jebakan: aturannya berlaku untuk \emph{bilangan pokok yang
sama} (atau \omterm{def:g8:powers:def}{eksponen} yang sama, untuk yang terakhir). Tidak ada aturan
yang menyederhanakan $2^3 \times 5^2$ --- hitung saja:
$8 \times 25 = 200$.
\end{example}

\section{Pangkat sepuluh}

\begin{proposition}[Pangkat sepuluh]\label{prop:g8:powers:ten}
Untuk $n \geq 1$: $10^n = 1\underbrace{0\dots0}_{n}$, dan
$10^{-n} = 0.\underbrace{0\dots0}_{n-1}1$. Aturan \omterm{def:g8:powers:def}{eksponennya}
berbunyi: mengalikan pangkat sepuluh berarti menjumlahkan \omterm{def:g8:powers:def}{eksponennya}.
\end{proposition}

\begin{example}\label{ex:g8:powers:ten}
$10^4 \times 10^3 = 10^7$; \quad
$\dfrac{10^2}{10^5} = 10^{-3} = 0.001$. Besaran yang besar dan yang
kecil menjadi terbaca:
\[
\text{satu miliar} = 10^9,
\qquad
\text{satu per sejuta} = 10^{-6}.
\]
Digabung dengan desimal: $3.2 \times 10^5 = 320\,000$ dan
$4.7 \times 10^{-3} = 0.0047$ (geserlah \omterm{def:g4:decimals:point}{titik desimalnya},
\cref{prop:g6:decimals:shift}). Pemakaian penulisan ini secara
sistematis --- \emph{notasi ilmiah} --- dikembangkan pada
\cref{ch:g9:fractions}.
\end{example}

\begin{example}[Orde besaran]\label{ex:g8:powers:magnitude}
Cahaya menempuh kira-kira $3 \times 10^8$ m/s; satu tahun punya
kira-kira $3.2 \times 10^7$ detik. Jadi satu tahun cahaya kira-kira
\[
3 \times 3.2 \times 10^{8+7} \approx 10 \times 10^{15} = 10^{16}
\text{ m}
\]
--- sepuluh juta miliar meter. Pangkat sepuluh membuat perhitungan
astronomi muat dalam satu baris.
\end{example}

\begin{method}[Menyederhanakan bentuk berpangkat]\label{met:g8:powers:simplify}
\begin{enumerate}
  \item Kelompokkan faktornya menurut bilangan pokoknya;
  \item terapkan aturan \omterm{def:g8:powers:def}{eksponennya} di dalam tiap bilangan pokoknya;
  \item hitunglah pangkat kecil yang tersisa, atau biarkan jawabannya
        berupa pangkat kalau nilainya besar.
\end{enumerate}
\end{method}

\begin{example}\label{ex:g8:powers:simplify}
\[
\frac{3^5 \times 3^{-2} \times 4^2}{3^2}
= \frac{3^{5 + (-2)}}{3^2} \times 16
= 3^{3 - 2} \times 16
= 3 \times 16 = 48 .
\]
\end{example}

\section{Latihan}

\begin{exercise}[$\star$]\label{exo:g8:powers:1}
Hitunglah:
\[
2^6, \qquad 3^3, \qquad 10^5, \qquad 1^{100}, \qquad 0.1^2, \qquad
6^0 .
\]
\end{exercise}

\begin{exercise}[$\star$]\label{exo:g8:powers:2}
Hitunglah:
\[
(-3)^2, \qquad -3^2, \qquad (-1)^{15}, \qquad (-5)^3, \qquad
\left(\tfrac{2}{3}\right)^2 .
\]
\end{exercise}

\begin{exercise}[$\star$]\label{exo:g8:powers:3}
Tulislah sebagai satu pangkat:
\[
7^4 \times 7^5, \qquad
\frac{2^9}{2^3}, \qquad
\left(10^3\right)^4, \qquad
5^6 \times 5^{-2}, \qquad
3^4 \times 7^4 .
\]
\end{exercise}

\begin{exercise}[$\star$]\label{exo:g8:powers:4}
Tulislah sebagai \omterm{def:g6:decimals:places}{bilangan desimal}: $10^{-2}$; \quad $4 \times 10^3$; \quad
$2.5 \times 10^{-4}$; \quad $10^0$.
\end{exercise}

\begin{exercise}[$\star$]\label{exo:g8:powers:5}
Tulislah dengan pangkat sepuluh: seratus ribu; sepersepuluh; sepuluh
miliar; $0.000\,001$.
\end{exercise}

\begin{exercise}[$\star$]\label{exo:g8:powers:6}
Sederhanakan, lalu hitunglah:
\[
\frac{10^7 \times 10^{-3}}{10^2},
\qquad
\frac{2^5 \times 2^4}{2^6},
\qquad
\left(2^2\right)^3 \times 2^{-4} .
\]
\end{exercise}

\begin{exercise}[$\star$]\label{exo:g8:powers:7}
Benar atau salah? Perbaikilah yang salah.
\[
2^3 \times 2^4 = 2^{12}; \qquad
5^2 + 5^3 = 5^5; \qquad
(3^2)^4 = 3^8; \qquad
10^3 \times 10^3 = 100^3 .
\]
\end{exercise}

\begin{exercise}[$\star\star$]\label{exo:g8:powers:8}
Sebuah kabar burung menyebar: pada hari 1, tiga orang mengetahuinya;
setiap hari, tiap orang yang mengetahuinya memberi tahu tiga orang
baru. Tulislah dengan pangkat banyaknya orang \emph{baru} yang diberi
tahu pada hari $4$, lalu hitunglah berapa orang yang mengetahui kabar
itu pada akhir hari 4 (termasuk ketiga orang pertamanya).
\end{exercise}

\begin{exercise}[$\star\star$]\label{exo:g8:powers:9}
Selembar kertas tebalnya $0.1$ mm, yaitu $10^{-4}$ m. Melipatnya
membuat tebalnya berlipat dua setiap kali.
\begin{enumerate}
  \item Nyatakan tebalnya sesudah $10$ lipatan sebagai \omterm{def:g2:mult:def}{hasil kali},
        lalu hitunglah dalam sentimeter ($2^{10} = 1024$).
  \item Sesudah $42$ lipatan tebalnya akan menjadi
        $2^{42} \times 10^{-4}$ m, dengan $2^{42} \approx 4.4 \times
        10^{12}$. Tunjukkan bahwa itu melampaui jarak Bumi--Bulan,
        yang kira-kira $3.8 \times 10^8$ m.
\end{enumerate}
\end{exercise}

\begin{exercise}[$\star\star$]\label{exo:g8:powers:10}
Urutkan dari yang terkecil ke yang terbesar, tanpa kalkulator:
\[
2^{10}, \qquad 10^3, \qquad 3^6, \qquad 5^4 .
\]
(Hitunglah masing-masing; $2^{10}$ dan $10^3$ adalah tetangga yang terkenal.)
\end{exercise}

\begin{exercise}[$\star\star\star$]\label{exo:g8:powers:11}
Mana yang lebih besar, $2^{100}$ atau $10^{30}$? Pakailah
$2^{10} = 1024 > 10^3$ untuk membandingkan
$2^{100} = \left(2^{10}\right)^{10}$ dengan
$\left(10^3\right)^{10}$.
\end{exercise}

\section{Soal: Papan catur dan pangkat dua}

\begin{problem}[{Soal akhir pekan --- \omterm{def:g1:addition:def}{jumlah} geometri
$1 + 2 + 4 + \dots + 2^{n-1} = 2^n - 1$, dari legenda yang terkenal
sampai bilangan biner}]
\label{pb:g8:powers:1}
Legendanya: sebagai hadiah karena menciptakan catur, si bijak Sissa
meminta kepada rajanya satu butir gandum pada petak pertama papannya,
dua pada petak kedua, empat pada petak ketiga --- berlipat dua dari
petak ke petak, sampai petak keenam puluh empat. Rajanya menertawakan
kesederhanaan itu. Soal ini menghitung apa yang dijanjikan rajanya,
dengan aturan \omterm{def:g8:powers:def}{eksponen} \cref{thm:g8:powers:rules}, lalu berakhir di
tempat yang diam-diam dituju kisahnya: bilangan biner di dalam setiap
komputer.

\medskip
\noindent\textbf{Bagian I --- Trik pelipatduaan.}
Untuk $n \geq 1$, misalkan $S_n$ banyaknya seluruh butir gandum pada
$n$ petak pertamanya.
\begin{enumerate}
  \item Nyatakan banyaknya butir pada petak $k$ sebagai pangkat
        $2$. Pangkat yang mana yang duduk pada petak $64$?
  \item Hitunglah $S_1$, $S_2$, $S_3$, $S_4$ dan $S_5$, lalu
        bandingkan masing-masing dengan pangkat $2$ yang terdekat.
        Dugalah rumus untuk $S_n$.
  \item \emph{Trik pelipatduaan}: tulislah \omterm{def:g1:addition:def}{jumlah} $S_n$ dan
        $2 \times S_n$ yang satu di bawah yang lain, kurangkan, lalu
        buktikan dugaanmu:
        \[
        S_n = 1 + 2 + 4 + \dots + 2^{n-1} = 2^n - 1 .
        \]
  \item Berapa butir gandum yang dijanjikan rajanya seluruhnya?
        Nyatakan jawabannya dengan pangkat $2$, lalu lengkapi
        catatan klasiknya: ``seluruh papannya memuat satu butir lebih
        sedikit daripada yang akan dimuat satu petak keenam puluh
        lima saja.''
  \item Tunjukkan bahwa separuh kedua papannya (petak $33$ sampai
        $64$) memuat \emph{tepat} $2^{32}$ kali sebanyak butir
        separuh pertamanya.
\end{enumerate}

\medskip
\noindent\textbf{Bagian II --- Seberapa besarkah $2^{64}$?}
Perbandingan $2^{10} = 1024 > 10^3$ pada
\cref{exo:g8:powers:11} adalah kunci semua taksiran di bawah ini.
\begin{enumerate}[resume]
  \item Tunjukkan bahwa
        $2^{64} = 2^4 \times \left(2^{10}\right)^6 > 1.6 \times
        10^{19}$.
  \item Satu butir gandum beratnya kira-kira $0.05$ g, yaitu
        $5 \times 10^{-2}$ g. Tunjukkan bahwa gandum yang dijanjikan
        itu beratnya lebih dari $8 \times 10^{17}$ g, lalu ubahlah
        menjadi ton ($1$ ton $= 10^6$ g).
  \item Seluruh dunia saat ini memanen kira-kira $8 \times 10^8$
        ton gandum tiap tahunnya. Sedikitnya berapa tahun panen
        dunia yang dijanjikan rajanya?
  \item Carilah bilangan bulat terkecil $n$ sehingga
        $2^n > 10^6$ --- yaitu berapa kali pelipatduaan yang
        diperlukan untuk melampaui satu juta. (Hitunglah $2^{19}$
        dan $2^{20}$ dengan persis, memakai $2^{10} = 1024$.)
  \item Seorang penipu menawarimu gaji sebulan: $1$ sen pada hari
        $1$, lalu dua kali bayaran hari sebelumnya pada tiap
        harinya. Pada hari yang mana bayaran \emph{harian}nya saja
        pertama kali melampaui satu juta euro ($10^8$ sen)?
        (Hitunglah $2^{26}$ dan $2^{27}$ dengan persis.)
\end{enumerate}

\medskip
\noindent\textbf{Bagian III --- Timbangan biner.}
Seorang pedagang punya lima anak timbangan: $1$, $2$, $4$, $8$ dan
$16$ gram, masing-masing satu. Ia menaruh beberapa di antaranya pada
satu piring neraca untuk menimbang barang di piring yang lain.
\begin{enumerate}[resume]
  \item Anak timbangan mana yang ditaruhnya untuk menimbang $21$ g?
        Untuk menimbang $27$ g?
  \item Jelaskan mengapa sasaran $16$ g atau lebih \emph{harus}
        memakai anak timbangan $16$ g, dan mengapa sasaran $15$ g
        atau kurang \emph{tidak boleh} memakainya. (Pertanyaan 3
        memberi tahu berapa paling banyak yang dapat dicapai anak
        timbangan $1, 2, 4, 8$.) Jelaskan mengapa penalaran yang
        sama berulang dengan anak timbangan terbesar berikutnya,
        pada setiap tahapnya.
  \item Simpulkan bahwa setiap sasaran bulat dari $1$ sampai $31$ g
        dapat ditimbang, dan \emph{dengan tepat satu cara}: setiap
        bilangan antara $1$ dan $31$ adalah \omterm{def:g1:addition:def}{jumlah} pangkat $2$ yang
        berbeda-beda dengan cara yang tunggal.
  \item Pedagangnya membeli anak timbangan keenam, seberat $32$ g.
        Sampai sasaran berapa ia kini dapat menimbang? Tulislah
        $45$ g sebagai \omterm{def:g1:addition:def}{jumlah} pangkat $2$ yang berbeda-beda.
  \item Di dalam komputer, bilangan ``64-bit'' disimpan pada $64$
        petak, masing-masing memuat $0$ atau $1$ --- petak $k$
        menyumbang $2^{k-1}$ ketika ia memuat $1$, seperti butir
        gandum pada legendanya. Dengan memakai Bagian I, jelaskan
        mengapa bilangan bulat yang dapat disimpan mesin seperti itu
        berjalan tepat dari $0$ sampai
        $2^{64} - 1$.

\end{enumerate}
\end{problem}
